\section{Introduction}

Vision transformers use patch tokens to represent image detail and
self-attention to integrate information across the image. Processing a fine
patch grid through every block, however, makes this global reasoning
expensive. Token merging reduces the sequence by combining patches rather
than simply removing them~\citep{tome2023}. A learned merger thus faces two
linked requirements: learn where information should flow, and physically
reduce the sequence seen by downstream blocks during training.

Differentiable merging does not automatically solve this execution problem.
DTEM learns dedicated merge embeddings through continuous grouping and
merging~\citep{dtem2024}. Its relaxed operators retain all token slots; for
end-to-end training, DTEM alternates soft updates of the merge embeddings
with hard, token-reducing updates of the backbone. Thus DTEM uses both
continuous policy learning and physical compression during training, but in
separate update paths. We ask whether one forward pass can propagate gradients
through soft routing while downstream blocks process an exact, physically
shorter sequence.

We study this question with \mn{}\footnote{An anonymized implementation is
provided in the supplementary material.}, a local-to-latent vision transformer. Six
local blocks encode the full patch grid, after which soft mass transport
learns how patch information moves between slots. A hard gather then selects
$K$ carriers from $N$ patches, so the subsequent latent blocks receive a
physically shorter tensor during both training and inference. A residual
cross-attention read lets those carriers access the full-grid local features
before global latent reasoning. The gather indices are discrete; a
straight-through weight connects the selected features to continuous routing
scores without claiming to differentiate the indices themselves.

The routing support follows the original two-dimensional patch grid. A donor
can send mass only to eligible receivers within radius $R$, while learned
similarities determine the weights and the resulting mass determines which
carriers survive the gather. Keeping these coordinates fixed gives the
routing support a direct interpretation on the image grid at every step.

Carrier mass also has a defined role after selection. Following proportional
attention in ToMe~\citep{tome2023}, \mn{} adds each carrier's $\log m$ to its
attention key logits. This per-key bias is rank one across queries; the
released operator reproduces it by adding one query/key dimension and using
the original softmax scale in an unmodified FlashAttention
call~\citep{flashattn2022}. We use an existing
threshold-based smooth top-$k$ relaxation~\citep{sujianlin2024topk} for the
continuous selection scores. Our contribution is the way these components
connect soft spatial transport to an exact-budget latent path, together with
the mass-bias implementation.

We evaluate the joint architecture on ImageNet-1K across carrier budgets,
input resolutions, and inference implementations. At 224 pixels, \mn{}
reaches 81.360\% top-1 while halving the patch sequence passed to its
latent stage; higher-resolution fine-tuning raises top-1 to 82.618\%
at 512 pixels. The default configuration reduces estimated arithmetic
cost by 37.0\% relative to dense DeiT-S/8. Separate synthetic benchmarks
show how that reduction translates into execution cost: compiled
batch-64 inference takes 27.24\,ms on an A100, close to the 27.22\,ms
measured for compiled ToMe. These evaluations connect the architectural
bottleneck to measurable classification and execution operating points.

Our contributions are:
\begin{itemize}
  \item a local-to-latent architecture that connects differentiable spatial
  mass transport to an exact-budget gather within one forward pass;
  \item an exact query/key reparameterization that retains the log-mass
  attention bias in an unmodified FlashAttention call; and
  \item an empirical characterization of carrier-budget and resolution
  changes on ImageNet-1K, together with arithmetic, latency, and memory
  measurements of the resulting implementations.
\end{itemize}
